\chapter{Deadlock Detection}
\label{protocol}

\section{Introduction}

We extend locks to detect \ac{tttl} deadlocks and report them as runtime exceptions.
The proof below establishes detection and absence of false reports in an abstract model with explicit implementation assumptions.

\section{Deadlock Detection Protocol}

We have modified the default implementation of Java's {\tt ReentrantLock} to allow efficient runtime detection of \ac{tttl} deadlocks. It works as follows:

\begin{enumerate}
\item Each lock has a pointer to a thread, the owner of the lock, or {\tt null} when no thread owns that lock.
\item Each lock has an integer to represent its current state: 0 means the lock is free and no thread owns it (the \emph{unlocked} state), 1 means there is a thread that owns the lock (the \emph{locked} state). For simplicity, we are only interested on these two states. Nonetheless, in the implementation of {\tt ReentrantLock}, each time a thread owner acquires the same lock, this state would be incremented, and decremented each time the thread releases it.
\item Each thread has a thread-local list of pointers to locks it currently owns.
\item Each lock has a waiting queue of threads that are waiting to acquire it. Whenever a thread tries to obtain a lock when it's already acquired, the thread will add itself to the waiting queue before parking. Upon the event of releasing the lock, the owner of that lock will look for the first thread in the waiting queue and unpark it.
\item When a thread wants to acquire a lock, it will swap the current state to \emph{locked} if the current state is \emph{unlocked} atomically.
\begin{enumerate}
\item If the thread fails, it must be because the lock is already owned by some other thread, then it will add itself on the waiting queue for that lock. Finally, the thread will park.
\item Otherwise, the thread will set itself as the current owner of that lock and also add this lock to its thread-local list of pointers of locks it owns.
\end{enumerate}
\item When a thread is about to release a lock, the current owner pointer of that lock is set to {\tt null} and that lock is also removed from the thread-local list of owned locks. Finally, the lock state is changed to \emph{unlocked}.
\item Before parking, a thread will check whether there is a deadlock. When the current thread is unable to acquire its desired lock, it must be because another thread already owns it. It is possible to know who is the owner of any lock, so the current thread identifies the owner of its desired lock as the conflicting thread. Then the current thread will search on each lock of its list of owned locks if the conflicting thread is waiting for it.
\begin{enumerate}
\item If positive, then we have a circular dependency (current thread is stuck waiting for its desired lock and the conflicting thread is stuck waiting for a lock the current thread owns) and thus a deadlock exception will be raised.
\item Otherwise, the thread parks.
\end{enumerate}
\end{enumerate}

The proof uses an abstract model with the following assumptions:
\begin{enumerate}
\item Locks are exclusive, and each thread's owned-lock list is accurate. A blocked acquisition retains its locks and remains queued until it succeeds or a deadlock is reported; timeout and interrupt cancellation are outside this model.
\item Queue insertion precedes the thread's deadlock check. Insertions and checks have a single order consistent with each thread's program order; checks see all earlier insertions that remain active.
\item Each check observes ownership and active queue membership in one consistent state. A thread completes this check before parking.
\end{enumerate}
These assumptions describe the abstract detector, not a verified Java memory-model implementation.

\subsection{Proof}

Let $A$ own $a$ and request $b$, while $B$ owns $b$ and requests $a$, with $A \ne B$ and $a \ne b$.
Write $E_A$ and $E_B$ for their queue insertions, and $C_A$ and $C_B$ for their detection checks.
The relation $\prec$ denotes the model's event order, so $E_A \prec C_A$ and $E_B \prec C_B$.

\begin{Lem}
Under the model assumptions, a \ac{tttl} cycle cannot leave both threads parked without a detection.
\end{Lem}
\begin{proof}
Suppose both threads park without reporting a deadlock.
Without loss of generality, let $A$ enqueue first. Then
\begin{equation}
E_A \prec E_B \prec C_B.
\end{equation}
At $C_B$, $A$ still owns $a$ and remains queued for $b$: neither acquisition can succeed while the other thread retains its lock.
Thus $B$, which owns $b$ and requests $a$, observes $A$ in $b$'s queue and reports the cycle before parking, a contradiction.
\end{proof}

\begin{Lem}
Under the model assumptions, every detection identifies a \ac{tttl} cycle.
\end{Lem}
\begin{proof}
Suppose $A$ detects a deadlock while requesting $b$.
Its check observes that $B$ owns $b$ and actively waits for a lock $a$ owned by $A$.
Because these observations describe one state, $A$ waits for $B$ and $B$ waits for $A$.
Neither can complete its acquisition and release its held lock without intervention: this is a \ac{tttl} cycle, even if a thread has not yet parked.
\end{proof}

The implementation must justify the observation assumptions.
Separate owner and queue reads are insufficient: $A$ could read $B$ as $b$'s owner, then $B$ could release $b$ and queue for $a$ before $A$ scans that queue.
The scan would then suggest a cycle although $b$ is free.
This sketch therefore leaves visibility, concurrent ownership changes, and cancellation as implementation proof obligations.

\section{Protocol Extension}

The base detector reports a cycle in at least one thread. The extension records the other participant so that it also receives an exception if it later acquires the awaited lock.
This requires the first thread to release that lock and the second to acquire it eventually; throwing an exception alone does not release a lock.
The extension uses the following steps:

\begin{enumerate}
\item Each lock has a list of tainted threads. This list should only be read or updated by the owner of that lock, allowing immunity from interference without any extra synchronization cost.
\item Once a deadlock is detected and the current thread is about to raise a deadlock exception, it already knows which thread is conflicting with itself and which lock that thread desires. The current thread (the owner of the desired lock) will add this conflicting thread to the tainted threads list for that lock. After that, the deadlock exception is raised.
\item When the conflicting thread is unparked and finally acquires its desired lock (it becomes the owner of that lock), then it is allowed to read the list of tainted threads. If this thread identifies itself in this list, then it must be because it was part of a deadlock before, so it removes its reference from the list and also raises a deadlock exception.
\item Every operation on the list of tainted threads of any lock (either reading or inserting values) should be followed up by some cleanup on all references to threads that are no longer running.
\end{enumerate}

Both threads may instead detect the cycle independently before either releases its lock. However, for this particular case, this change introduces a different problem: dangling references.
If each thread adds their conflicting thread to the lists of tainted threads of the locks they own, 
but none of them is able to acquire their respective desired locks (as in \emph{item 3}),
both threads will leave their references behind for others to cleanup (as in \emph{item 4}).
We minimize this issue by cleaning these references as soon as any thread acquires the lock.

\section{Implementation}

The full version of the modified OpenJDK {\tt ReentrantLock} that implements this algorithm is available in our code repository~\citep{repo}.
In our initial commit, we copied the original {\tt ReentrantLock} source code and related lock classes from Java SE Development Kit 8 and kept it mostly without changes, and then afterwards we landed commits that actually implement the concepts we will discuss in this section. 

We will focus on what changes were done on the original implementation presenting pseudo-code to illustrate how they work (more details on our repository's commits history).

\medskip
\noindent
{\it Keeping ownedLocks list updated}
\begin{verbatim}
Lock currentLock;
DEFINE_PER_LOCK(Thread, owner);

DEFINE_PER_THREAD(List<Lock>, ownedLocks);

setExclusiveOwner(Thread t) {
  currentLock.owner = t;
  if (currentLock.owner == null) {
    unregisterOwnedLock();
  } else {
    registerOwnedLock();
  }
}

registerOwnedLock() {
  ownedLocks.push(currentLock);
}

unregisterOwnedLock() {
  ownedLocks.remove(currentLock);
}
\end{verbatim}

Following the first part of the protocol, we must keep a list of locks each thread owns.
This list is thread-local: the main advantage is that we automatically get freedom from interference as each thread will only manage its own list.
In original implementation, {\tt setExclusiveOwner} is used whenever a thread successfully acquire a lock or release it.
On our version, we modified it to intercept it and also update the list of owned locks accordingly as follows:
whenever the call resets the owner it means there was a release, so we unregister that lock and removes it from the list of owned locks of the current thread;
futhermore, we do the opposite when the owner is not null, as it means that the thread has owned that lock and it should update its own list of owned locks to add that particular entry.

\medskip
\noindent
{\it Detecting deadlock and throw exception}
\begin{verbatim}
park() {
  Thread other = currentLock.owner;
  if (isAnyOwnedLockDesiredBy(other)) {
    clearOwnedLocksByCurrentThread();
    throw new DeadlockException();
  }
  LockSupport.park(this);
}

isAnyOwnedLockDesiredBy(Thread t) {
  forEach(lock : ownedLocks) {
    if (lock.queue.contains(t)) {
      return True;
    }
  }
  return False;
}

clearOwnedLocksByCurrentThread() {
  ownedLocks.clear();
}
\end{verbatim}

When a thread attempts to acquire a lock and this lock is already owned, we deploy the deadlock check right before the thread get parked.
It starts by checking which thread is the owner of this particular lock.
Then the current thread checks on its list of owned locks whether any of them has the owner of the desired lock waiting for it.
If positive, then we have a circular wait: this thread wants another lock, but the owner of that lock is waiting for a lock this thread owns.
In that case, we must throw a deadlock exception.

The changes described previously were part the first version of the algorithm and we minimally changed {\tt ReentrantLock} to support it.
That version was used on our usability evaluation (Section~\ref{sec:usab}).
The base detector only guarantees a report in at least one participant under the proof assumptions.

The extension notifies the second participant if it subsequently acquires the awaited lock.
We will present again some functions to illustrate how they changed in this last version.

\medskip
\noindent
{\it Changes on ReentrantLock: throw deadlock exception on both threads}
\begin{verbatim}
DEFINE_PER_LOCK(List<Thread>, taintedThreads);

park() {
  Thread other = currentLock.owner;
  List<Lock> desiredLocks;
  desiredLocks = getMyOwnedLocksDesiredBy(other);
  if (!desiredLocks.empty()) {
    forEach(lock : desiredLocks) {
      lock.taintedThreads.add(other);
    }
    clearOwnedLocksByCurrentThread();
    throw new DeadlockException();
  }
  LockSupport.park(this);
}

getMyOwnedLocksDesiredBy(Thread t) {
  List<Lock> desiredLocks = [];
  forEach(lock : ownedLocks) {
    if (lock.queue.contain(t)) {
      desiredLocks.add(lock);
    }
  }
  return desiredLocks;
}

void setExclusiveOwner(Thread t) {
  currentLock.owner = t;
  if (currentLock.owner == null) {
    unregisterOwnedLock();
  } else {
    registerOwnedLock();
    if (taintedThreads.contains(t)) {
      clearOwnedLocksByCurrentThread();
      throw new DeadlockException();
    }
  }
}

\end{verbatim}

We expanded the algorithm by adding a list of threads on each lock object containing threads that just went into a deadlock and should also throw exception as soon as possible.
Note although this particular list is not thread-local, we guaranteed freedom from interference,
as we applied a policy on that list of tainted threads to only permit read/write operations for the owner of the lock.

When a deadlock was detected on the previous version, the following case was possible:
one thread throws an exception and releases all its locks,
then the second thread would finally acquire its desired lock
which should be free after the first thread released its locks.

However, it is important to note that when the first thread is about to throw an exception,
the second thread is already stuck waiting for the first thread.

Also, the first thread was still the owner of the current lock object allowing it to modify the list of tainted threads inside that lock.
Then it could update list of tainted threads by adding the second thread on it.
Finally, by the time the second thread acquires this lock which should be free now, it should also throw an exception.

This extension can retain {\tt Thread} references inside tainted threads lists for each lock.
For example, both threads may detect the deadlock and throw exceptions before either acquires the awaited lock.
In that case, they would have added the opposite thread inside their own lock's taintedThreads list, but afterwards
both threads would stop and none of them would attempt to acquire the opposite lock again.
We've minimized the effect of this leak by adding a policy to remove non-active thread
references from tainted threads list every time any update is done on it.
In practice we expect other threads would eventually clean up these thread references.
We also provide code that implements this last policy on our repository~\cite{repo}.
